\documentclass[11pt]{article}
\usepackage[margin=1.1in]{geometry}
\usepackage{amsmath,amssymb,amsthm}
\usepackage{algorithm}
\usepackage[noend]{algpseudocode}
\usepackage{booktabs}
\usepackage{placeins}
\usepackage[hidelinks]{hyperref}
\usepackage[capitalise,nameinlink]{cleveref}

\newtheorem{theorem}{Theorem}
\newtheorem{lemma}[theorem]{Lemma}
\newtheorem{corollary}[theorem]{Corollary}
\newtheorem{proposition}[theorem]{Proposition}
\theoremstyle{definition}
\newtheorem{definition}[theorem]{Definition}
\newtheorem{remark}[theorem]{Remark}
\newtheorem{example}[theorem]{Example}

\newcommand{\Nm}{N^{-}}
\newcommand{\vs}{\operatorname{vs}}
\newcommand{\pw}{\operatorname{pw}}
\newcommand{\Rel}{\mathcal{R}}
\newcommand{\Gr}{\mathcal{G}}
\newcommand{\cost}{w}
\newcommand{\bott}{b}
\newcommand{\set}[1]{\{#1\}}

\title{One-pass bottleneck search for linear width parameters:\\
exact pathwidth and cutwidth, optimal orderings,\\
and secondary costs}
\author{Victor Miller\\ {\small IDA Center for Communications Research}}
\date{\today}

\begin{document}
\maketitle

\begin{abstract}
Kobayashi~\cite{Kob15} gave a very simple dynamic program that decides
whether a directed (or undirected) graph $D$ has pathwidth at most $k$ in
time $3^{\tau(D)} n^{O(1)}$. Here $\tau(D)$ is the vertex cover number of the
underlying undirected graph. The dynamic program runs over \emph{relevant}
vertex sets, which are sets closed under a natural ``subsumption''
expansion. Used as a decision procedure, the algorithm must be called
repeatedly to find the pathwidth. Each call with $k$ above the true value is
expensive, because the number of $k$-feasible relevant sets grows quickly
with $k$.

We observe that the cost of a step of Kobayashi's dynamic program does not
depend on $k$. As a result, the vertex separation number is the value of a
minimum-bottleneck path in a directed acyclic graph on relevant sets. A
bucketed Dijkstra-type search computes it in a single pass. That pass
discovers only relevant sets of bottleneck at most $\pw(D)$, so it never
does more work than one call of the decision procedure at $k=\pw(D)$, and in
practice it does considerably less. The same search, with one parent pointer
per set, yields an explicit vertex ordering of width $\pw(D)$. We prove this,
and we point out a subtlety: within one expansion step the vertices must be
emitted in the order in which the expansion subsumes them. We describe the
data structures of Python and C implementations and report timings.

Two further results concern the use of the search. First, a greedy rule from
the branch-and-bound algorithm of Coudert, Mazauric and Nisse can be added to
the expansion; this leaves every guarantee intact and reduces the number of
sets 2--6 times. Second, when the last level is completed, the relevant sets
form a DAG whose paths are pathwidth-optimal orderings: about $10^{26}$ of
them for the square-sum graph on $\{1,\dots,57\}$. Any secondary cost that
is local to a step can be minimised over all of these exactly, including
costs with a multiplicative carry such as those used to tune frontier
dynamic programs. On our test graphs the C implementation is about 6--20
times faster than the Cython branch and bound in Sage; that implementation
exceeded 120\,s on four instances with its default limit on stored prefixes.

Finally, we ask which of these ideas apply to other width parameters. Both
greedy rules used for pathwidth are instances of one lemma: for a
submodular prefix cost, a vertex whose addition does not increase the cost
can be added at once. For vertex separation, the two rules together are
exactly the non-increasing additions. The lemma also applies to cutwidth.
Combined with the vertex-cover normal form of Cygan, Lokshtanov, Pilipczuk,
Pilipczuk and Saurabh, the same one-pass search computes cutwidth exactly
over at most $2^{\tau}$ states, and in a Python prototype it gives the
cutwidth of the square-sum graph on $\{1,\dots,60\}$. For treewidth, the
search applies to elimination orderings but its greedy rules are weak.
Bandwidth does not fit the framework, because its cost depends on positions
and not on prefix sets.
\end{abstract}

\section{Introduction}

Let $D=(V,A)$ be a directed graph with $n=|V|$ vertices. For $X\subseteq V$
write $\Nm(X)=\bigl(\bigcup_{x\in X}\Nm(x)\bigr)\setminus X$, where $\Nm(x)$
is the set of in-neighbours of~$x$. A linear ordering $\sigma=(v_1,\dots,v_n)$
of $V$ has \emph{width}
\[
  \cost(\sigma)=\max_{1\le i\le n}\bigl|\Nm(\set{v_1,\dots,v_i})\bigr|,
\]
and the \emph{vertex separation number} $\vs(D)$ is the minimum width over
all orderings. By Yang and Cao~\cite{YC08}, $\vs(D)$ equals the directed
pathwidth $\pw(D)$. For an undirected graph $G$, replace each edge by a pair
of antiparallel arcs; then $\Nm$ is the ordinary open neighbourhood, and
$\vs(G)=\pw(G)$ by Kinnersley~\cite{Kin92}. Everything below applies to both
cases, and we write $N(\cdot)$ for $\Nm(\cdot)$ when no confusion is possible.

Computing pathwidth is NP-hard~\cite{KF79}. Exact methods fall roughly into
three families.

\emph{Mathematical programming.} Coudert~\cite{Cou16} gave compact integer
linear programming formulations for several linear ordering problems,
pathwidth among them. His variables record whether a vertex belongs to a
prefix of the ordering, rather than its exact position. An implementation
accompanies the paper~\cite{CouCode}. Mallach~\cite{Mal18} studied
formulations based on linear-ordering variables. The present author
rewrote Coudert's vertex separation model as a MaxSAT instance
(\cref{sec:maxsat}). On the author's test graphs, the MaxSAT version was frequently hundreds of times
faster than the MIP solved with Gurobi~\cite{Gurobi}.

\emph{Branch and bound.} Coudert, Mazauric and Nisse~\cite{CMN14} explore
layouts depth first. They combine three devices: reduction rules that
contract parts of the graph, greedy steps that extend a prefix without loss
(their Lemma~6), and a table of prefix \emph{sets} already explored (their
Lemma~7). Their implementation is part of Sage~\cite{Sage}. Two of these
devices reappear below. The first case of their greedy step is exactly
Kobayashi's subsumption, with arcs reversed, and their second case can be
added to the expansion of \cref{sec:alg} (\cref{sec:greedy2}). Their prefix
table is a memo on the same states as the set dynamic program.

\emph{Dynamic programming over vertex sets.} The classical recurrence of
Bodlaender et al.~\cite{BFKKT12} runs over all $2^n$ prefix sets. Pruning
tricks reduce this substantially~\cite{SV09,KKKTT12,Tam11}. In particular,
Kobayashi~\cite{Kob15} showed that one may restrict the dynamic program to
\emph{relevant} sets, of which there are at most $3^{\tau(D)}$. This gives
the $3^{\tau(D)}n^{O(1)}$ bound of Chapelle et al.~\cite{CLTV13}, but with a
much simpler algorithm, and one that extends to digraphs.

Kobayashi's Algorithm~1 is a decision procedure: given $k$, it answers
whether $\vs(D)\le k$. Two practical issues arise in using it.

\begin{enumerate}
\item \textbf{Locating the optimum.} Binary search over $k$ works, but every
probe with $k>\pw(D)$ is expensive: fewer sets are pruned, and the number
$|\Rel_k|$ of relevant sets reachable under bound $k$ grows rapidly with $k$
(\cref{tab:overshoot}). Probes below $\pw(D)$ are wasted work as well.
\item \textbf{Producing a witness.} The decision procedure does not output an
ordering of width $\pw(D)$.
\end{enumerate}

\paragraph{Other widths.} The approach solved these problems for pathwidth
far faster than the mathematical programming models, so it is natural to
ask which other width parameters it fits. Cutwidth, treewidth and
bandwidth are the classical candidates. Cutwidth parameterized by the
vertex cover number is already known to be solvable in $2^{\tau}n^{O(1)}$
time~\cite{CLPPS14}, by an argument closely related to the greedy rules
used here. For treewidth, the most effective exact solvers use
positive-instance driven dynamic programming over minimal separators and
potential maximal cliques~\cite{Tam19}.

\paragraph{Contributions.}
\begin{itemize}
\item We show that the cost of each transition of Kobayashi's dynamic
program is independent of $k$ (\cref{lem:mono}). Consequently $\vs(D)$ is a
minimum-bottleneck path value in a DAG on relevant sets (\cref{thm:bott}).
\item We give a one-pass, bucketed bottleneck search
(\cref{alg:bucket}). It computes $\vs(D)$ without any width parameter. It
discovers every relevant set of bottleneck $<\pw(D)$, and only some of those
of bottleneck exactly $\pw(D)$. So its work lies between that of
Algorithm~1 at $k=\pw(D)-1$ and at $k=\pw(D)$ (\cref{thm:work}). It retains
the $3^{\tau(D)}n^{O(1)}$ bound.
\item We show how to recover an optimal ordering from parent pointers
(\cref{thm:order}), and we identify the necessary ordering constraint inside
an expansion step (\cref{ex:order}).
\item We show that the second greedy rule of~\cite{CMN14} can be added to
the expansion without affecting any of these guarantees
(\cref{sec:greedy2}).
\item Completing the last level of the search yields a DAG whose paths are
pathwidth-optimal orderings. We show how to minimise any step-local secondary
cost over all of them exactly, including costs with a multiplicative carry
(\cref{sec:orders}).
\item We describe compact data structures (\cref{sec:ds}) and report
timings for Python and C implementations, together with comparisons against
a MaxSAT model and the branch and bound of~\cite{CMN14} (\cref{sec:exp}).
\item We prove a greedy lemma for any submodular prefix cost
(\cref{lem:submod}). For vertex separation, its hypothesis is exactly the
union of Kobayashi's subsumption and the second rule of~\cite{CMN14}
(\cref{prop:vsrules}).
\item Using the lemma and the normal form of~\cite{CLPPS14}, the one-pass
search computes cutwidth exactly over at most $2^{\tau}$ states
(\cref{sec:cutwidth}). We also explain why treewidth gains little and
bandwidth does not fit (\cref{sec:twbw}).
\end{itemize}

\section{Kobayashi's algorithm}\label{sec:kob}

We recall the definitions of~\cite{Kob15}. A sequence $\sigma$ of distinct
vertices is \emph{$k$-feasible} if every prefix $\sigma_i$ satisfies
$|N(V(\sigma_i))|\le k$. A set $U\subseteq V$ is $k$-feasible if
$U=V(\sigma)$ for some $k$-feasible $\sigma$, and \emph{strongly
$k$-feasible} if moreover $\sigma$ is a prefix of a $k$-feasible ordering of
all of $V$. Thus $\vs(D)\le k$ if and only if $V$ is $k$-feasible.

\begin{definition}
Let $U\subseteq V$. A vertex $x\in V\setminus U$ is \emph{subsumed} by $U$
if $\Nm(x)\subseteq U\cup N(U)$. The \emph{expansion} $U^{*}$ is obtained by
repeatedly adding subsumed vertices to $U$ until none remain. $U$ is
\emph{relevant} if $U^{*}=U$.
\end{definition}

\begin{lemma}[{\cite[Lemma~1]{Kob15}}]\label{lem:kob1}
If $U$ is $k$-feasible (resp.\ strongly $k$-feasible), so is $U^*$.
\end{lemma}

\begin{algorithm}[t]
\caption{Kobayashi's decision procedure~\cite[Algorithm~1]{Kob15}.}\label{alg:kob}
\begin{algorithmic}[1]
\For{$i=0,\dots,n$} $R_i\gets\emptyset$ \EndFor
\State add $\emptyset^{*}$ to $R_{|\emptyset^{*}|}$
\For{$i=0,\dots,n-1$}
  \ForAll{$U\in R_i$}
    \ForAll{$v\in V\setminus U$ with $|N(U\cup\set{v})|\le k$}
      \State add $(U\cup\set v)^{*}$ to $R_{|(U\cup\set v)^*|}$
    \EndFor
  \EndFor
\EndFor
\State \Return ``yes'' if $R_n\ne\emptyset$, else ``no''
\end{algorithmic}
\end{algorithm}

\begin{remark}
In the published pseudocode, line~2 reads ``$R_i\gets\{\emptyset\}$''.
This must mean $R_i\gets\emptyset$, the empty collection: read literally,
$R_n$ would always be nonempty.
\end{remark}

Kobayashi proves that \cref{alg:kob} is correct~\cite[Lemma~2]{Kob15}. He
also proves that there is an injection from relevant sets into ordered
tripartitions of a minimum vertex cover~\cite[Lemma~3]{Kob15}. Hence at most
$3^{\tau(D)}$ sets are ever stored.

\section{The cost of a step is independent of \texorpdfstring{$k$}{k}}\label{sec:mono}

\begin{lemma}\label{lem:mono}
If $x\in V\setminus U$ is subsumed by $U$, then $N(U\cup\set x)\subseteq N(U)$.
\end{lemma}
\begin{proof}
$N(U\cup\set x)=\bigl(N(U)\cup \Nm(x)\bigr)\setminus(U\cup\set x)$. Since
$\Nm(x)\subseteq U\cup N(U)$, we have
$\Nm(x)\setminus(U\cup\set x)\subseteq N(U)$.
\end{proof}

This is the observation at the start of the proof of~\cite[Lemma~1]{Kob15};
the point is that it involves no $k$. Iterating it gives the following.

\begin{corollary}\label{cor:exp}
Let $T\subseteq V$, and let $x_1,\dots,x_m$ be the vertices of
$T^*\setminus T$ in the order an expansion procedure adds them. That is,
each $x_j$ is subsumed by $T_{j-1}=T\cup\set{x_1,\dots,x_{j-1}}$. Then
$N(T_j)\subseteq N(T)$ for all $0\le j\le m$. In particular
$|N(T_j)|\le|N(T)|$.
\end{corollary}

\begin{definition}[Relevant-set DAG]
Let $\Gr$ be the directed graph whose vertices are the relevant sets. For
each relevant $U\ne V$ and each $x\in V\setminus U$, it has an arc
$U\to(U\cup\set x)^*$ labelled $x$, with cost
\[\cost(U,x)=|N(U\cup\set x)|.\]
Every arc strictly enlarges the set, so $\Gr$ is acyclic. The cost of a path
is the maximum cost of its arcs; the empty path has cost $0$. For a relevant
set $S$, let $\bott(S)$ be the minimum cost of a path from $\emptyset^*$ to
$S$, or $\infty$ if there is none.
\end{definition}

\begin{definition}[Ordering of a path]\label{def:pathorder}
Let $P=(S_0\xrightarrow{x_1}S_1\to\dots\xrightarrow{x_m}S_m)$ be a path in
$\Gr$ with $S_0=\emptyset^*$. Its \emph{ordering} $\sigma(P)$ is the
concatenation of the following segments:
\begin{itemize}
\item the vertices of $\emptyset^*$, in the order the expansion procedure
adds them;
\item for each $i=1,\dots,m$, the vertex $x_i$, followed by the vertices of
$S_i\setminus(S_{i-1}\cup\set{x_i})$ in the order the expansion procedure,
started from $S_{i-1}\cup\set{x_i}$, adds them.
\end{itemize}
\end{definition}

\begin{lemma}\label{lem:pathwidth}
$\sigma(P)$ is a sequence of distinct vertices with $V(\sigma(P))=S_m$, and
every prefix $\pi$ of it satisfies
$|N(V(\pi))|\le\max_i\cost(S_{i-1},x_i)$. Equality holds for the prefix
ending at $x_i$, where $i$ attains the maximum.
\end{lemma}
\begin{proof}
First, prefixes inside the first segment: by \cref{cor:exp} with
$T=\emptyset$, they have empty boundary. Next, consider a prefix ending
within segment $i$. Either it equals $S_{i-1}\cup\set{x_i}$, whose boundary
has size exactly $\cost(S_{i-1},x_i)$, or it is $T_j$ for
$T=S_{i-1}\cup\set{x_i}$ in the notation of \cref{cor:exp}, whose boundary
is contained in $N(T)$.
\end{proof}

\begin{theorem}\label{thm:bott}
$\bott(V)=\vs(D)$.
\end{theorem}
\begin{proof}
$(\ge)$: if $\bott(V)<\infty$, a path of cost $\bott(V)$ yields, by
\cref{lem:pathwidth}, an ordering of $V$ of width $\bott(V)$.

$(\le)$: let $k=\vs(D)$, so $V$ is $k$-feasible. The proof
of~\cite[Lemma~2]{Kob15} gives the following. Start from
$\emptyset^*$, which is strongly $k$-feasible by \cref{lem:kob1}. From a
strongly $k$-feasible relevant $U\ne V$ there is $v$ with $U\cup\set v$
strongly $k$-feasible; in particular $|N(U\cup\set v)|\le k$, the
boundary of the last prefix. Then $(U\cup\set v)^*$ is again strongly
$k$-feasible and relevant. Repeating this traces a path from $\emptyset^*$
to $V$ in $\Gr$ with all costs $\le k$.
\end{proof}

\begin{remark}\label{rem:Rk}
For each $k$, the collection $\Rel_k=\bigcup_i R_i$ that \cref{alg:kob}
computes with bound $k$ is exactly $\set{S:\bott(S)\le k}$. Indeed,
\cref{alg:kob} processes sets in order of size, and arcs strictly increase
size, so induction on $|S|$ applies. In particular
$\Rel_0\subseteq\Rel_1\subseteq\cdots$.
\end{remark}

Two cheap consequences are already useful if one keeps the decision
procedure.

\begin{proposition}\label{prop:next}
Suppose \cref{alg:kob} with bound $k$ answers ``no''. Let
$k'=\min\set{\cost(U,x): U\in\Rel_k,\ \cost(U,x)>k}$, the smallest cost
rejected at line~5. Then $\Rel_j=\Rel_k$ for $k\le j<k'$, and so
$\vs(D)\ge k'$.
\end{proposition}
\begin{proof}
For $j$ in this range no new arc becomes admissible, so $\Rel_j=\Rel_k$
by \cref{rem:Rk}. Since $V\notin\Rel_k$, we get $\vs(D)=\bott(V)>j$ for
every such $j$.
\end{proof}

\begin{proposition}\label{prop:loose}
Run \cref{alg:kob} with any $k\ge\vs(D)$. Whenever $(U\cup\set v)^*=S$ is
generated, update
$f(S)\gets\min\bigl(f(S),\max(f(U),\cost(U,v))\bigr)$, with
$f(\emptyset^*)=0$. Then $f(V)=\vs(D)$ at the end.
\end{proposition}
\begin{proof}
All predecessors of $S$ are strictly smaller, so they are final when $S$'s
size class is processed. Hence $f(S)=\bott(S)$ for every $S\in\Rel_k$.
Moreover $\Rel_k$ contains every set on an optimal path to $V$, since those
sets have $\bott\le\vs(D)\le k$. Finally apply \cref{thm:bott}.
\end{proof}

\Cref{prop:loose} gives the exact value from a single overestimate, but it
pays the full cost $|\Rel_k|$. The next section avoids that cost.

\section{A one-pass bucketed bottleneck search}\label{sec:alg}

Path cost in $\Gr$ is a maximum, so it is monotone along paths. A
Dijkstra-type search~\cite{Dij59} therefore applies to the
minimum-bottleneck (``maximum capacity'') path problem~\cite{Pol60}. The
costs are integers in $[0,n-1]$, so a priority queue is unnecessary: we
process levels $c=0,1,2,\dots$ in order.

\begin{algorithm}[t]
\caption{One-pass bottleneck search with parent pointers.}\label{alg:bucket}
\begin{algorithmic}[1]
\State $S_0\gets\emptyset^*$; \ $\mathit{par}[S_0]\gets\bot$; \ $c\gets 0$
\State $B_0,\dots,B_n\gets\emptyset$ \Comment{level buckets of discovered sets}
\State $W\gets\emptyset$ \Comment{work stack of pairs $(S,x)$ for the current level}
\Statex
\Procedure{Scan}{$S,\ell,c$}\Comment{release arcs of cost in $(\ell,c]$; reschedule $S$}
  \State $c'\gets\infty$
  \ForAll{$x\in V\setminus S$}
    \State $w\gets|N(S\cup\set x)|$
    \If{$\ell<w\le c$} push $(S,x)$ on $W$
    \ElsIf{$w>c$} $c'\gets\min(c',w)$ \EndIf
  \EndFor
  \If{$c'<\infty$} append $S$ to $B_{c'}$ \EndIf
\EndProcedure
\Statex
\If{$S_0=V$} \Return $(0,\ \text{expansion order of }\emptyset)$ \EndIf
\State \Call{Scan}{$S_0,-1,0$}
\Loop
  \While{$W\neq\emptyset$}
    \State pop $(P,x)$ from $W$; \ $T\gets(P\cup\set x)^*$
    \If{$T\in\mathrm{dom}(\mathit{par})$} \textbf{continue} \EndIf
    \State $\mathit{par}[T]\gets(P,x)$
    \If{$T=V$} \Return $(c,\ \textsc{Extract}(\mathit{par}))$ \EndIf
    \State \Call{Scan}{$T,-1,c$}
  \EndWhile
  \State $c\gets c+1$
  \ForAll{$S\in B_c$} \Call{Scan}{$S,c-1,c$} \EndFor
\EndLoop
\end{algorithmic}
\end{algorithm}

Every set discovered at some level is placed in exactly one bucket, namely
that of its cheapest not-yet-released outgoing arc. When the search reaches
that level, the set is rescanned, the arcs of exactly that cost are released,
and the set is rescheduled. Consequently each arc $(S,x)$ is pushed onto
$W$ exactly once, at level $\max(c_S,\cost(S,x))$, where $c_S$ is the level
at which $S$ was discovered.

\begin{theorem}[Correctness]\label{thm:correct}
\Cref{alg:bucket} discovers every relevant set $S$ with $\bott(S)<\vs(D)$,
and each set it discovers at level $c$ satisfies $\bott(S)=c$. It returns
$c=\vs(D)$.
\end{theorem}
\begin{proof}
We prove by induction on $c$ the following claim: when level $c$ ends (that
is, $W$ is empty and the search is about to move to $c+1$), every $S$ with
$\bott(S)\le c$ has been discovered, and every set discovered so far has
$\bott$ equal to its discovery level.

\emph{Discovery levels are correct.} Suppose $T$ is discovered at level $c$
via a pair $(P,x)$. That pair was pushed at level
$\max(c_P,\cost(P,x))\le c$. By induction $\bott(P)=c_P$, so
$\bott(T)\le c$. On the other hand, $T$ was not discovered by the end of
level $c-1$, so by the induction hypothesis $\bott(T)>c-1$. Hence
$\bott(T)=c$.

\emph{Completeness.} Let $\bott(S)\le c$, and take a path of cost $\le c$
from $S_0$ to $S$. Every set on it has $\bott\le c$. Induct along the
path. If a set $S'$ on it has been discovered, at level $c_{S'}\le c$, then
the next arc $(S',x)$ has cost $\le c$, so it is pushed by the end of
level $c$: at discovery if $\cost\le c_{S'}$, or at the rescan of level
$\cost(S',x)$ otherwise, since $S'$ is rescheduled at each successive
larger cost. The level-$c$ loop runs until $W$ is empty, so the head of the
arc is discovered during level $c$ or earlier.

The first time $V$ is discovered, at level $c$, we therefore have
$c=\bott(V)=\vs(D)$ by \cref{thm:bott}. Termination is clear, because
$\bott(V)\le n-1$.
\end{proof}

\begin{theorem}[Optimal ordering]\label{thm:order}
Chase the parent pointers from $V$ back to $\emptyset^*$. This gives a path
$P$ in $\Gr$ of cost exactly $\vs(D)$, and $\sigma(P)$ from
\cref{def:pathorder} is an ordering of $V$ of width $\vs(D)$.
\end{theorem}
\begin{proof}
If $\mathit{par}[T]=(P,x)$, then $(P,x)$ was pushed at a level $\le c_T$,
so $\cost(P,x)\le c_T$, and $P$ was discovered at level $c_P\le c_T$. By
induction, the path of parent pointers to any $T$ has cost $\le c_T$. For
$T=V$ this is $\le\vs(D)$. \Cref{lem:pathwidth} gives an ordering of that
width, and \cref{thm:bott} gives the reverse inequality.
\end{proof}

The proof shows that \emph{any} recorded parent is acceptable: each arc
popped at level $c$ has cost $\le c$, and its tail was discovered at a level
$\le c$. Hence the first parent found is optimal, and no comparison between
parents is ever needed. (In the setting of \cref{prop:loose}, by contrast,
one must keep the parent that attains the minimum of $f$.)

\begin{example}[Order within an expansion matters]\label{ex:order}
In \cref{def:pathorder} the expansion vertices of a step cannot be emitted
in arbitrary order. Consider an undirected graph with vertices
$a,b,x,w_1,w_2,w_3$ and edges $ab$, $ax$, and $w_ib$, $w_ix$ for $i=1,2,3$.
Take $T=\set a$, so $N(T)=\set{b,x}$. Each $w_i$ is subsumed by $T$, since
$N(w_i)=\set{b,x}\subseteq T\cup N(T)$. The vertex $x$ is not, since
$N(x)=\set{a,w_1,w_2,w_3}$. Once the $w_i$ have been added, $x$ (and $b$)
become subsumed. If $x$ were placed immediately after $a$, the prefix
$\set{a,x}$ would have boundary $\set{b,w_1,w_2,w_3}$ of size $4>2=|N(T)|$.
Emitting vertices in the order the expansion adds them, one at a time with
each subsumed by the current set, avoids this. Kobayashi's own definition
adds vertices in rounds; any order within a round is fine, but the order of
the rounds must be kept.
\end{example}

\begin{theorem}[Work]\label{thm:work}
Let $\mathcal{D}$ be the family of sets discovered by \cref{alg:bucket}.
Then
\[
  \Rel_{\vs(D)-1}\ \subseteq\ \mathcal{D}\ \subseteq\ \Rel_{\vs(D)},
\]
where $\Rel_{-1}=\emptyset$. Each $S\in\mathcal{D}$ is scanned at most $n$
times, and each arc is pushed at most once. Hence the running time is
$O(|\Rel_{\vs(D)}|\cdot n^{O(1)})=3^{\tau(D)}n^{O(1)}$.
\end{theorem}
\begin{proof}
The inclusions follow from \cref{thm:correct,rem:Rk}. A set is rescanned
only at levels equal to distinct costs of its outgoing arcs, of which there
are at most $n$. The bound $|\Rel_k|\le 3^{\tau(D)}$ is~\cite[Lemma~3]{Kob15}.
\end{proof}

\paragraph{Comparison with search over $k$.}
By \cref{thm:work}, the one-pass search costs at most one run of
\cref{alg:kob} at the optimal $k$. Binary search over
$k\in[\ell,u]$ costs $\sum_{\text{probes}}|\Rel_{k}|$, and that sum is
dominated by its largest probe, which in general exceeds $\vs(D)$. Even
linear search upward from a lower bound costs
$\sum_{k=\ell}^{\vs(D)}|\Rel_k|$ and repeats all the work below the optimum
at every probe. In our experiments $\mathcal{D}$ is close to the lower end
of the sandwich (\cref{tab:sandwich}). The reason is that $W$ is a LIFO
stack: at level $\vs(D)$ the search proceeds depth-first and usually reaches
$V$ after discovering only a small fraction of
$\Rel_{\vs(D)}\setminus\Rel_{\vs(D)-1}$.

\paragraph{Decision version.}
To decide $\vs(D)\le k$, stop when $c$ would exceed $k$. By
\cref{thm:work}, this explores a subset of $\Rel_k$, and it returns early
with a witness ordering when $\vs(D)\le k$.

\subsection{A stronger expansion}\label{sec:greedy2}

Coudert, Mazauric and Nisse~\cite[Lemma~6]{CMN14} prove that a prefix can be
extended greedily, without loss, by a vertex $v$ that satisfies either of
two conditions. They state the lemma for the out-neighbour convention;
applied to the reversed digraph it reads as follows in ours.

\begin{lemma}[{\cite[Lemma~6]{CMN14}}]\label{lem:cmn6}
Let $U\subsetneq V$ and $v\in V\setminus U$. Suppose that either
\textup{(i)} $\Nm(v)\subseteq U\cup N(U)$, or \textup{(ii)} $v\in N(U)$ and
$|\Nm(v)\setminus(U\cup N(U))|=1$. If $U$ is strongly $k$-feasible, so is
$U\cup\set v$.
\end{lemma}

Condition~(i) is subsumption. Under~(ii), with $w$ the unique in-neighbour
of $v$ outside $U\cup N(U)$, we have $N(U\cup\set v)=(N(U)\setminus\set
v)\cup\set w$, so the boundary keeps its size. Define $U^{**}$ by adding
vertices one at a time, each satisfying~(i) or~(ii) with respect to the
current set, by a fixed deterministic rule, until neither applies.

\begin{proposition}\label{prop:greedy2}
All results of \cref{sec:mono,sec:alg} remain true when $U^*$ is replaced by
$U^{**}$ throughout, with the family $\Rel_k$ of \cref{rem:Rk} replaced by
the sets generated under $U^{**}$. In particular \cref{alg:bucket} still
returns $\vs(D)$ together with an ordering of that width, and every set it
stores is relevant, so the bound $3^{\tau(D)}n^{O(1)}$ still holds.
\end{proposition}
\begin{proof}
\Cref{cor:exp} is replaced by the remark that each vertex added under~(i)
or~(ii) leaves the boundary size unchanged or makes the boundary smaller.
Hence every prefix inside a step has boundary size at most $\cost(U,x)$,
and \cref{lem:pathwidth} holds. In the proof of \cref{thm:bott}, the
appeal to \cref{lem:kob1} is replaced by \cref{lem:kob1,lem:cmn6}. The
proofs of \cref{thm:correct,thm:order,thm:work} use only these two facts,
together with the determinism of the expansion. Finally, $U^{**}$ is closed
under~(i), so it is relevant.
\end{proof}

Conditions (i) and (ii) together say exactly that adding $v$ does not
increase $|N(U)|$ (\cref{prop:vsrules}). \Cref{lem:cmn6} is therefore a
special case of the general greedy lemma of \cref{sec:submod}.

$U^{**}$ may depend on the order in which the rules are applied: (ii) is
not monotone. The arguments need only a fixed procedure. On our test graphs
the stronger expansion reduces the number of stored sets 2--6 times
(\cref{tab:bab}). It does reduce the family of orderings in
\cref{sec:orders} much more sharply, so there we use $U^*$.

\section{All relevant optimal orderings, and secondary costs}\label{sec:orders}

A dynamic program that is driven by a vertex ordering, such as a
frontier dynamic program for counting Hamiltonian paths, often depends on
far more than the width. Many orderings attain $\pw(D)$, and their running
times can differ by an order of magnitude. This section shows that the
search already contains a compact representation of a large family of
optimal orderings, and that step-local secondary costs can be minimised
exactly over that family.

\begin{definition}\label{def:gamma}
Let $k=\vs(D)$. The \emph{complete DAG} $\Gamma$ has as nodes $V$ and the
relevant sets $S$ with $\bott(S)\le k$ from which $V$ can be reached by arcs
of cost $\le k$; its arcs are the arcs of $\Gr$ of cost $\le k$ between
nodes.
\end{definition}

\begin{proposition}\label{prop:gamma}
The map $P\mapsto\sigma(P)$ of \cref{def:pathorder} is injective on paths
of $\Gamma$ from $\emptyset^*$ to $V$, and every $\sigma(P)$ has width
exactly $\vs(D)$. $\Gamma$ is obtained from \cref{alg:bucket} by continuing
level $c=\vs(D)$ until $W$ is empty, instead of stopping at $V$, and then
making one backward pass over the stored sets in order of decreasing size.
\end{proposition}
\begin{proof}
By \cref{lem:pathwidth}, the width of $\sigma(P)$ is the largest arc cost
on $P$, which is at most $k$, so it equals $k$. Two different paths first
differ at a node $S$ by their arc label $x$, which is the next vertex of
the ordering. The proof of \cref{thm:correct}, applied with $c=k$, shows
that completing level $k$ discovers every relevant set of bottleneck at most
$k$, that is, $\Rel_k$. The backward pass keeps the sets with an arc of cost
$\le k$ into a kept set or into $V$; the arcs out of $S$ are recomputed
from $S$ and need not be stored.
\end{proof}

The paths of $\Gamma$ do not give every optimal ordering: they are the
orderings that take subsumed vertices eagerly. In the convention of a
left-to-right frontier dynamic program, the description is cleaner. The
frontier of a prefix $T$ is $F(T)=\set{u\in T:N(u)\not\subseteq T}$, and by
\cref{lem:rev}, $\sigma$ has width $k$ exactly when $\tau=\sigma^R$ has
$\max_i|F(T_i)|=k$. For an undirected graph, a set $S$ is relevant exactly
when every vertex of $T=V\setminus S$ has a neighbour in $T\setminus F(T)$,
that is, a neighbour that has already left the frontier. (Indeed, $x\in T$
is subsumed by $S$ iff $N(x)\subseteq S\cup N(S)=(V\setminus T)\cup F(T)$.)
So the orderings of $\Gamma$, read as frontier orderings, are those
in which this holds at every block boundary.

\paragraph{Step-local costs.}
Fix the frontier direction. A step adds a vertex $v$ to a prefix $T$, and
we call a cost \emph{step-local} if each step contributes through
quantities that depend only on $(T,v)$. Examples are the new frontier size
$|F(T\cup\set v)|$, the number of edges from $v$ to vertices still in the
frontier, the number of frontier vertices eliminated, the number of edges
inside the frontier, the component sizes of the subgraph it induces, and
the step index $|T|$. We allow a multiplicative carry $p$, as in proxy
models of the state count of a dynamic program. A step contributes a triple
$(a,e,\kappa)$ of nonnegative reals that depends only on $(T,v)$, and acts
by
\[
\text{cost}\gets\text{cost}+a+p\,e,\qquad p\gets\max(p\,e\,\kappa,\,1),
\]
starting from $p=1$. Additive costs are the case $e=0$.

\begin{proposition}\label{prop:pareto}
Let $J(T,p)$ be the least cost of completing the prefix $T$ from carry $p$.
Then $J(T,\cdot)$ is nondecreasing. Consequently, a dynamic program over
$\Gamma$, run in the frontier direction, that keeps at each node only the
labels $(\text{cost},p)$ not dominated in both coordinates by another label
there, returns the exact minimum over all orderings of $\Gamma$. For
additive costs the front has a single label, and keeping the $k$ smallest
costs per node gives the $k$ best orderings.
\end{proposition}
\begin{proof}
$J(V,p)=0$, and
$J(T,p)=\min_{v}\bigl[a+p\,e+J\bigl(T\cup\set v,\max(p\,e\,\kappa,1)\bigr)\bigr]$.
By induction on $|V\setminus T|$, each term is nondecreasing in $p$
($e,\kappa\ge0$), and hence so is the minimum. If $c'\le c$ and $p'\le p$,
then $c'+J(T,p')\le c+J(T,p)$, so a dominated label is never needed. In
$\Gamma$ a step is a block; its triples are applied in sequence, and the
frontier direction traverses $\Gamma$ backwards, from $V$ (the empty
frontier prefix) to $\emptyset^*$.
\end{proof}

The same dynamic program, run over all prefixes $T$ with every
$|F(T_i)|\le k$ rather than over $\Gamma$, gives the minimum over
\emph{all} optimal orderings. It is the reference against which we measure
the restriction to $\Gamma$ (\cref{tab:gammaloss}); that family is much
larger (\cref{sec:exp}).

\section{Other width parameters}\label{sec:widths}

\Cref{alg:bucket} uses only three properties of vertex separation.
\begin{enumerate}
\item \emph{Set-determined steps.} The cost $c(S,v)\ge0$ of placing $v$
after a prefix set $S$ depends only on $S$ and $v$.
\item \emph{A min--max objective.} The width of an ordering is the largest
step cost along it.
\item \emph{Greedy closure.} Some extensions of a prefix are safe whatever
the width bound, and the cost of an arc is the largest step cost in its
block.
\end{enumerate}
Under~(1) and~(2) the width is a minimum-bottleneck path value in the DAG of
prefix sets. The proofs of \cref{thm:correct,thm:order} then apply verbatim,
with arcs $S\to\mathrm{cl}(S\cup\set v)$ and arc cost the largest step cost
in the block. (When the closure can raise that cost, as for treewidth below,
the block must be computed when the arc is scanned.) They use only the monotonicity of a maximum along a path and
the determinism of the closure $\mathrm{cl}$. Property~(3) is what makes the
search small. The bound $3^{\tau}$ of \cref{thm:work} is specific to
Kobayashi's closure. This section asks which width parameters have the
three properties.

\subsection{A greedy lemma for submodular prefix costs}\label{sec:submod}

Call $f\colon 2^V\to\mathbb{R}$ a \emph{prefix cost} when $c(S,v)=f(S\cup\set
v)$. Then the width of an ordering is $\max_i f(S_i)$ over its prefixes.
Write $\Delta(S,v)=f(S\cup\set v)-f(S)$. Recall that $f$ is
\emph{submodular} exactly when $\Delta(\cdot,v)$ is antitone:
$\Delta(X,v)\ge\Delta(Y,v)$ for $X\subseteq Y\not\ni v$.

\begin{lemma}\label{lem:submod}
Let $f$ be a submodular prefix cost. Suppose $S$ is a prefix of an ordering
of value $M$ (so $f(S)\le M$), and $v\notin S$ satisfies $\Delta(S,v)\le0$.
Then $S\cup\set v$ is a prefix of an ordering of value at most~$M$.
\end{lemma}
\begin{proof}
Let the ordering continue after $S$ with $a_1,\dots,a_m,v,\dots$, and move
$v$ to the front of the continuation. With $A_j=\set{a_1,\dots,a_j}$, each
new prefix $S\cup A_j\cup\set v$ ($0\le j\le m$) satisfies
\[
f(S\cup A_j\cup\set v)=f(S\cup A_j)+\Delta(S\cup A_j,v)\le
f(S\cup A_j)+\Delta(S,v)\le f(S\cup A_j)\le M.
\]
Every other prefix is unchanged.
\end{proof}

For vertex separation, $f(S)=|\Nm(S)|$ is submodular. Indeed
$S\cup\Nm(S)$ is a coverage function of $S$, so it is submodular, and $|S|$
is modular.

\begin{proposition}\label{prop:vsrules}
For $f(S)=|\Nm(S)|$ and $v\notin U$, we have $\Delta(U,v)\le0$ exactly when
condition~\textup{(i)} or~\textup{(ii)} of \cref{lem:cmn6} holds.
Consequently the strong-feasibility half of Kobayashi's \cref{lem:kob1}
and \cref{lem:cmn6} both follow from \cref{lem:submod}, applied one vertex at a time, and the closure $U^{**}$
of \cref{sec:greedy2} consists exactly of the non-increasing additions.
\end{proposition}
\begin{proof}
Write $w=|\Nm(v)\setminus(U\cup N(U))|$. Then $N(U\cup\set v)=(N(U)\setminus\set
v)\cup(\Nm(v)\setminus(U\cup N(U)))$, so $\Delta(U,v)=w-[v\in N(U)]$. This
is $\le0$ exactly when $w=0$, which is~(i), or when $v\in N(U)$ and $w=1$,
which is~(ii).
\end{proof}

We checked \cref{prop:vsrules} by brute force on 73\,080 pairs $(U,v)$ from
random graphs and digraphs. We also checked \cref{lem:submod} directly, for
vertex separation and for cutwidth, on 200 random graphs and digraphs with
$n\le9$.

\begin{remark}
The mirror image of \cref{lem:submod} is false. One might expect that a
vertex placed while $\Delta>0$ can be delayed until $\Delta\le0$, but the
prefix immediately before its new position is not dominated. On the path
$2-1-0-3$ with $f$ the cut size, the optimal ordering $(2,1,0,3)$ has width
$1$. Delaying $2$ until after $1$ creates the prefix $\set1$ of cut~$2$.
Delaying is allowed only while $\Delta\ge0$ along the way~\cite{CLPPS14}.
\end{remark}

Other submodular prefix costs include the cut size (\cref{sec:cutwidth})
and the cut-rank function $\rho(S)=\operatorname{rank}_{\mathrm{GF}(2)}A[S,V\setminus
S]$. The latter defines linear rank-width, which we have not tested.

\subsection{Cutwidth}\label{sec:cutwidth}

The cutwidth $\mathrm{cw}(G)$ of an undirected graph is the minimum over
orderings of $\max_i|E(S_i,V\setminus S_i)|$. The cut size $f(S)$ is
submodular. Here $\Delta(S,v)=\deg v-2|N(v)\cap S|$, which is the
\emph{rank} of~\cite{CLPPS14}, and the first part of their Lemma~1 is
\cref{lem:submod} for this~$f$. So \cref{alg:bucket} applies with:
\begin{itemize}
\item arc cost $f(S\cup\set x)$;
\item closure: repeatedly add any vertex of rank $\le0$, which never raises
the cost of the arc.
\end{itemize}

\paragraph{Vertex-cover restriction.} Let $C$ be any vertex cover, and
$I=V\setminus C$. Cygan et al.~\cite[Lemma~2]{CLPPS14} show that some
optimal ordering places every vertex of $I$ at a rank of $0$ or $\pm1$, at
a position determined by $S\cap C$. Odd-degree vertices of $I$ that are
balanced, meaning rank $+1$ now and $-1$ once their next neighbour arrives,
may go just before that neighbour, and only their number matters. We
therefore choose arc labels only from~$C$. For $x\in C\setminus S$, let
$Y(S,x)$ be the set of odd-degree vertices of $I\setminus S$ adjacent to
$x$ that have rank $+1$ at~$S$. The arc first places $i$ of them, then $x$,
then applies the closure, which absorbs the rest of $Y(S,x)$. Its cost is
\[
f(S)+\min_{0\le i\le|Y(S,x)|}\max\bigl(i,\ \Delta(S,x)-i\bigr),
\]
because the two cuts involved are $f(S)+i$ and $f(S)+\Delta(S,x)-i$.

\begin{proposition}\label{prop:cw2tau}
With the vertex-cover restriction, every stored set $S$ satisfies $S\cap
I=\set{u\in I:\Delta(S\cap C,u)\le0}$. Hence at most $2^{|C|}$ sets are
stored, and the search runs in $2^{|C|}n^{O(1)}$ time.
\end{proposition}
\begin{proof}[Sketch]
A vertex $u\in I$ has all its neighbours in $C$, so its rank depends only
on $S\cap C$, and it is antitone in $S\cap C$. Vertices of $I$ enter only
through the closure, or as members of $Y(S,x)$, whose rank is $-1$ once $x$
is placed. Once a vertex of $I$ has rank $\le0$ it keeps rank $\le0$, and
the closure adds every vertex of rank $\le0$.
\end{proof}

This matches the bound of~\cite{CLPPS14}, whose dynamic program runs over
pairs (subset of $C$, last vertex of $C$). Our version adds the closure on
$C$ as well, and the one-pass search settles only states of bottleneck at
most $\mathrm{cw}(G)$. We have not written out a full proof that the
closure on $C$ is compatible with the normal form of~\cite{CLPPS14}. Instead
we checked the whole procedure against a brute-force $2^n$ dynamic program:
1\,520 runs on 560 random and bipartite graphs with $n\le13$, over up to
three variants (no rules, closure only, closure plus cover), with no
disagreement.
A first attempt that delayed vertices of $I$ naively until rank
$\le0$ (see the remark in \cref{sec:submod}) was one too high on 75 of 400
graphs, which is how the subtlety was found.

\Cref{tab:cw} reports a Python prototype on the square-sum graphs. The
cover is a minimum vertex cover, found by simple branching. Every reported
width is attained by an ordering that was re-checked, and is consistent with
$\pw\le\mathrm{cw}$. For $n\ge45$ the optimality of the value rests on the
argument above and the small-graph tests.

\begin{table}[ht]
\centering
\small
\begin{tabular}{rrrrrrrr}
\toprule
$n$ & $\tau$ & $\vs$ & $\mathrm{cw}$ & closure + cover & time (s) & closure only & no rules\\
\midrule
25 & 13 & 4 & 6 & 141 & 0.0 & 372 & 7\,871\\
30 & 16 & 5 & 6 & 191 & 0.0 & 785 & 3\,589\\
35 & 19 & 5 & 9 & 1\,292 & 0.2 & 6\,519 & 31\,306\\
40 & 21 & 7 & 11 & 3\,465 & 0.7 & 29\,085 & 199\,503\\
45 & 25 & 9 & 13 & 12\,065 & 2.8 & & \\
50 & 29 & 9 & 14 & 22\,345 & 6.8 & & \\
55 & 32 & 10 & 17 & 94\,068 & 37 & & \\
57 & 33 & 11 & 20 & 393\,746 & 168 & & \\
60 & 34 & 13 & 22 & 771\,807 & 359 & & \\
\bottomrule
\end{tabular}
\caption{Exact cutwidth of the square-sum graphs $Q_n$ with the one-pass
search (Python prototype, one core). The columns ``closure + cover'',
``closure only'' and ``no rules'' give the number of settled states.
$\tau$ is the vertex cover number. The stored states are far fewer than
$2^{\tau}$, because the search never goes above the optimum.}\label{tab:cw}
\end{table}

The data structures of \cref{sec:ds} carry over directly to a C version:
only the cost function, the rank closure and the cover-restricted arcs
change.

\subsection{Treewidth and bandwidth}\label{sec:twbw}

\paragraph{Treewidth.} The treewidth is the minimum over elimination
orderings of $\max_i|Q(S_{i-1},v_i)|$~\cite{BFKKT12}. Here $Q(S,v)$ is the
set of vertices outside $S\cup\set v$ reachable from $v$ through~$S$. This
cost is set-determined, so properties~(1) and~(2) hold and the one-pass
search applies. Eliminating simplicial vertices of the current elimination
graph is a safe closure: such a vertex has degree at most the treewidth of
what remains. Unlike the closures above, it can raise the cost of its arc,
so the arc cost is the maximum over the block. There is no submodular prefix
function here, so \cref{lem:submod} does not apply. In our prototype (which
agrees with brute force) the simplicial rule only reduced the settled states
on $Q_{20}$ from 61\,963 to 20\,788. The elimination-ordering dynamic
program is in any case far weaker than Tamaki's positive-instance driven
dynamic program over minimal separators and potential maximal
cliques~\cite{Tam19}. The ingredient that does transfer to that setting is
replacing the outer loop over $k$ by a single bottleneck pass. Because
positive-instance driven programs combine pairs of blocks, this needs
Knuth's generalisation of Dijkstra's algorithm to superior
functions~\cite{Knu77}. We expect a modest constant-factor gain, since the
run at the optimal $k$ usually dominates, but have not measured it. A
$3^{\tau}n^{O(1)}$ algorithm for treewidth is due to Chapelle et
al.~\cite{CLTV13}.

\paragraph{Bandwidth.} The bandwidth of an ordering is $\max_{uv\in
E}|\operatorname{pos}u-\operatorname{pos}v|$. This depends on positions
within the prefix, not only on the prefix set. Two prefixes with equal sets
but different orders have different futures, because each placed vertex $u$
forces its unplaced neighbours to land by position
$\operatorname{pos}u+b$. The classical dynamic programs~\cite{Sax80,GS84}
therefore keep the last $b$ placed vertices in order, so the state space
itself depends on the bound $b$. Property~(1) fails, and so do the exchange
arguments, because moving a vertex shifts every vertex it passes. What does
transfer is a lower bound. An ordering of bandwidth $b$ has every prefix
boundary of size at most $b$, so $\pw(G)\le\mathrm{bw}(G)$, and the
pathwidth computed here can seed an incremental SAT or ILP
search~\cite{BBNNPR13}. Exact exponential algorithms are also
known~\cite{CP12}.

\section{Data structures}\label{sec:ds}

\paragraph{Bit sets.}
Vertices are numbered $0,\dots,n-1$, and a vertex set is a bit mask. In C
this is one \texttt{uint64\_t}, so $n\le 64$. In Python it is an arbitrary-precision
\texttt{int}. The graph is stored as the array $\mathit{nin}[x]$ of
in-neighbour masks; for an undirected graph these are the neighbourhood
masks. Every discovered set $S$ is carried together with
\[
  \nu(S)=\bigcup_{u\in S}\mathit{nin}[u],
\]
so that $N(S)=\nu(S)\mathbin{\&}\lnot S$. Then the cost of a candidate arc
is a constant number of word operations:
\[
  \cost(S,x)=\operatorname{popcount}\bigl((\nu(S)\mid\mathit{nin}[x])\mathbin{\&}\lnot(S\mid 2^x)\bigr),
\]
and $\nu(S\cup\set x)=\nu(S)\mid\mathit{nin}[x]$. We iterate over the
members of $V\setminus S$ using count-trailing-zeros and clearing the lowest
set bit.

\paragraph{Expansion.}
To compute $T^*$ we maintain the closed mask $C=T\mid\nu(T)$, which is
$T\cup N(T)$. A vertex $x\notin T$ is subsumed exactly when
$\mathit{nin}[x]\mathbin{\&}\lnot C=0$. We sweep $V\setminus T$, add every
subsumed vertex immediately (updating $T$, $\nu$ and $C$), and repeat until a
sweep adds nothing. This one-at-a-time procedure is exactly what
\cref{cor:exp} requires. It is deterministic, so the same expansion order
can be recomputed during extraction instead of being stored. Each sweep costs
$O(n)$ word operations, and there are at most $|T^*\setminus T|+1$ sweeps.

\paragraph{Visited table and parent pointers.}
Discovered sets are kept in a hash table mapping $S\mapsto(P,x)$. In C this
is open addressing with linear probing, a power-of-two capacity, load factor
$\le 1/2$, and the 64-bit finaliser of MurmurHash3 as the mixing function.
The key is stored as $S+1$, so that $0$ can mark an empty slot. This cannot
overflow, because the full set $V$ is never inserted: its parent link is
kept in a separate variable, since its discovery ends the search. A slot holds $S$, $P$ and $x$, which is 24 bytes with
padding. In Python the table is a \texttt{dict}. We store neither $\nu(S)$
nor the expansion order; both are recomputed in $O(n^2)$ word operations per
step during extraction, which runs once.

\paragraph{Pre-expansion test.}
Before computing $(P\cup\set x)^*$ we look up $T_0=P\cup\set x$ itself. If
$T_0$ is already in the table, it is relevant, so $T_0^*=T_0$ and it has
been discovered. The expansion can then be skipped. This avoids the most
expensive operation for a substantial fraction of popped pairs.

\paragraph{Work stack and level buckets.}
The work stack $W$ holds triples $(S,\nu(S),x)$. A naive bucketed Dijkstra
would place every rejected arc $(S,x)$ in the bucket of its cost, costing up
to $n$ entries per set. Instead, each discovered set appears in exactly one
bucket at a time, as a pair $(S,\nu(S))$: the bucket of its cheapest arc not
yet released (procedure \textsc{Scan}). When level $c$ is reached, $S$ is
rescanned, and its arcs of cost exactly $c$ are released. This trades $O(n)$
recomputation per rescan, which is cheap with bit-parallel cost evaluation,
for an $n$-fold reduction in bucket memory. Within a level the processing
order does not affect correctness (\cref{thm:correct}). LIFO order is used
because it reaches $V$ quickly at the final level and gives good cache
locality.

\paragraph{Memory.}
On the largest test graph (\cref{tab:time}), the C implementation's peak
resident memory was about 390 bytes per discovered set, against about
1.2\,KB per set for Python. The C figure includes the hash table at load
$\le1/2$, the buckets, and the work stack. Since $|\mathcal{D}|$ can grow
like $3^{\tau(D)}$, memory rather than time is usually the binding
constraint.

\paragraph{Larger graphs and further engineering.}
A second C program, \texttt{kpw2}, implements four extensions, and
\texttt{doc/kpw2\_design.md} in the code repository gives the details and
measurements.
\begin{enumerate}
\item \emph{More than 64 vertices.} Sets are arrays of $W$ words, fixed at
compile time, so $n\le64W$. Hash slots carry an explicit occupancy byte, so
no key value is reserved.
\item \emph{A restricted expansion sweep.} A vertex can become subsumed
only when $U\cup N(U)$ gains one of its in-neighbours. Rule-1 expansion
therefore runs as a worklist over the out-neighbours of the newly closed
vertices. The closure is the least fixpoint of a monotone operator, so the
sets, and all counts, are identical to those of the full sweep. On our
graphs the full sweep took 14--18\% of the time, which bounds the possible
saving.
\item \emph{Compact parents.} With an option, a slot stores only the key,
the arc label $x$ and $|P|$, which makes the table 40\% smaller. During
extraction, a predecessor is recovered by one table scan per step. The scan
looks for a stored $P\subseteq T\setminus\set x$ of the recorded size with
$\cost(P,x)\le\vs(D)$ whose expansion with $x$ is $T$. Any such $P$ gives a
valid optimal chain, and the recorded parent is one of them. We first tried
undoing the expansion by peeling vertices. That search branched over every
interior vertex and was abandoned.
\item \emph{Threads.} By \cref{thm:correct}, the order of work within a
level is irrelevant, so threads share each level. The hash table is
sharded with spin locks, each thread drains its own stack depth first and
steals half of another's when empty, and an atomic count of pending pairs
ends the level.
\end{enumerate}
A further change cut memory substantially. The stack holds one item per
set, carrying the mask of arc labels still to try, instead of one entry per
label. Taking the highest label of the top item reproduces the order of
\cref{alg:bucket} exactly. With one thread and $W=1$, \texttt{kpw2}
matches \texttt{kpw} in time and uses about a third of its memory (for
example, 57\,MB against 201\,MB on $Q_{63}$). With $W=2$ it costs about
$1.5\times$ in time. Four threads gave speedups of 2--2.5 on a 4-core
machine; the search is dominated by cache misses in the hash table. With
\texttt{kpw2}, the square-sum graphs $Q_n$ for odd $n$ from 65 to 75 take under a
second each with the stronger expansion. Their widths are
$\vs(Q_n)=14$ for $n=65,67,\dots,73$ and $\vs(Q_{75})=15$, and
$Q_{65}$, $Q_{67}$ and $Q_{75}$ were checked against the Python
implementation.

\section{The MaxSAT baseline}\label{sec:maxsat}

For comparison we used a MaxSAT version of Coudert's MIP model for vertex
separation~\cite{Cou16,CouCode}. We state it precisely here, because it
measures the boundary on the opposite side of the cut from \cref{sec:kob},
and the two conventions are related by reversing the ordering.

\paragraph{Variables.} For $v\in V$ and $1\le t\le n$:
\begin{itemize}
\item $y_{v,t}$, meaning ``$v$ is among the first $t$ vertices'';
\item $u_{v,t}$, meaning ``$v$ is counted in the boundary at time $t$'';
\item $z_t$, the unary digits of the objective.
\end{itemize}

\paragraph{Hard clauses.}
\begin{align}
&\lnot y_{v,t}\lor y_{v,t+1} && v\in V,\ 1\le t<n, \label{eq:mono}\\
&\textstyle\sum_{v} y_{v,t}=t && 1\le t\le n, \label{eq:card}\\
&\lnot y_{v,t}\lor y_{w,t}\lor u_{v,t} && (v,w)\in A,\ 1\le t\le n, \label{eq:u}\\
&z_t\lor\lnot z_{t+1} && 1\le t<n, \label{eq:zmono}\\
&\textstyle\sum_{v} u_{v,t}+\sum_{s=1}^{n}\lnot z_s\le n && 1\le t\le n. \label{eq:zbound}
\end{align}
The cardinality constraints \eqref{eq:card} and \eqref{eq:zbound} are
encoded with totalizers. For an undirected graph, each edge gives both arcs.

\paragraph{Soft clauses.} $\lnot z_t$ with weight $1$, for $1\le t\le n$.

Coudert's MIP also has variables $x_{v,t}$ for the closed neighbourhood of
the prefix, and counts $x_{v,t}-y_{v,t}$, which is the boundary
\emph{outside} the prefix. The MaxSAT version eliminates $x$. It counts
instead the vertices \emph{inside} the prefix that have an out-neighbour
outside it; this is the convention of Kinnersley's $V_L$~\cite{Kin92}.
Since $u$ only enters an upper bound that is being minimised, clause
\eqref{eq:u} needs to force $u_{v,t}$ in one direction only.

For an ordering $\sigma$ with prefix sets $S_t$, write
\[
L_t(\sigma)=\set{v\in S_t:\ N^{+}(v)\not\subseteq S_t},
\]
where $N^{+}(v)$ is the set of out-neighbours of $v$.

\begin{lemma}[Reversal]\label{lem:rev}
Let $\sigma^{R}$ be the reverse of $\sigma$. Then
$L_t(\sigma)=\Nm\bigl(V(\sigma^{R}_{n-t})\bigr)$ for $0\le t\le n$. Hence
$\max_t|L_t(\sigma)|=\cost(\sigma^{R})$.
\end{lemma}
\begin{proof}
$V(\sigma^{R}_{n-t})=V\setminus S_t$. The in-neighbours of $V\setminus S_t$
that lie outside it are exactly the vertices of $S_t$ with an arc into
$V\setminus S_t$.
\end{proof}

\begin{proposition}\label{prop:maxsat}
The optimal cost of the MaxSAT instance is $\vs(D)$. Given an optimal model,
let $f(v)=\min\set{t:y_{v,t}}$, and order the vertices by increasing $f$.
The reverse of that ordering has width $\vs(D)$ in the convention of
\cref{sec:kob}.
\end{proposition}
\begin{proof}
By \eqref{eq:mono} and \eqref{eq:card}, the sets
$S_t=\set{v:y_{v,t}}$ form a chain with $|S_t|=t$, so $f$ is a bijection
onto $[n]$ and defines an ordering $\sigma$. Clause \eqref{eq:u} forces
$u_{v,t}$ for every $v\in L_t(\sigma)$. Let $Z$ be the number of true
$z_t$. By \eqref{eq:zmono} these are $z_1,\dots,z_Z$, so \eqref{eq:zbound}
reads $\sum_v u_{v,t}\le Z$, and therefore $Z\ge\max_t|L_t(\sigma)|$.
Conversely, setting $u_{v,t}$ true exactly on $L_t(\sigma)$ and
$Z=\max_t|L_t(\sigma)|$ satisfies every hard clause. So the optimum is
$\min_\sigma\max_t|L_t(\sigma)|$, which equals $\min_\sigma
\cost(\sigma^{R})=\vs(D)$ by \cref{lem:rev}.
\end{proof}

We checked \cref{prop:maxsat} against \cref{alg:bucket} on 400 random graphs
and digraphs with $n\le 9$. Both the optimal value and the width of the
reversed ordering agreed in every case.

The instance has $2n^2+n$ primary variables and $n|A|$ ternary clauses. Its
$2n$ totalizers contribute $O(n^3)$ clauses in all. We solved it with the
core-guided solver RC2~\cite{RC2} in PySAT~\cite{PySAT}, using the
CaDiCaL~1.9.5 backend. RC2 lowers the bound on $Z$ through a sequence of
unsatisfiable cores. In that respect it, like \cref{alg:bucket}, approaches
the optimum from below. Each of its oracle calls, however, is a SAT problem
over all $n$ time steps.

\section{Experiments}\label{sec:exp}

We tested the Python implementation against a brute-force $2^n$ dynamic
program on 3000 random graphs and digraphs with $n\le 10$. In every case the
returned value agreed, and the extracted ordering had width exactly
$\vs(D)$. The C port reports the same numbers of discovered sets, popped
pairs and rescans as the Python implementation on all the graphs below. Both
programs recompute the width of the extracted ordering directly and check
it against the returned value. All timings below are on one core of an
Intel Xeon at 2.1\,GHz, using CPython~3.13 and \texttt{gcc -O3
-march=native}~13.3. The graphs $g_1,g_2,g_3$ are Erd\H{o}s--R\'enyi graphs
$G(n,p)$ with $(n,p)=(40,0.12)$, $(60,0.06)$ and $(45,0.15)$. The graph
$Q_{63}$ is the \emph{square-sum graph} on $\set{1,\dots,63}$, in which
$i\sim j$ exactly when $i+j$ is a perfect square; it has 135 edges.

\begin{table}[ht]
\centering
\begin{tabular}{lrrrrrr}
\toprule
graph & $n$ & $\vs$ & $|\mathcal{D}|$ & Python (s) & C (s) & ratio\\
\midrule
$g_1$ & 40 & 11 & 4\,022 & 0.22 & 0.009 & $\approx 24$\\
$g_3$ & 45 & 15 & 19\,608 & 1.45 & 0.026 & $\approx 56$\\
$g_2$ & 60 & 11 & 732\,821 & 66--70 & 1.04 & $\approx 65$\\
$Q_{63}$ & 63 & 14 & 364\,589 & 25.6 & 0.43 & $\approx 59$\\
\bottomrule
\end{tabular}
\caption{One-pass search (\cref{alg:bucket}). The C time for $g_1$ is
dominated by process start-up. Peak memory on $g_2$: Python 854\,MB, C
288\,MB. On the author's own machine, $Q_{63}$ took 73.6\,s in Python and
0.965\,s in C, a ratio of about 76.}\label{tab:time}
\end{table}

\begin{table}[ht]
\centering
\begin{tabular}{lrrrr}
\toprule
graph & $|\Rel_{\vs-1}|$ & $|\mathcal{D}|$ & $|\Rel_{\vs}|$ & sets found at level $\vs$\\
\midrule
$g_1$ & 3\,959 & 4\,022 & 9\,500 & 63\\
$g_3$ & 18\,755 & 19\,608 & 44\,952 & 853\\
\bottomrule
\end{tabular}
\caption{The sandwich of \cref{thm:work}. Here $|\Rel_k|$ is the number of
sets generated by \cref{alg:kob} with bound $k$.}\label{tab:sandwich}
\end{table}

\begin{table}[ht]
\centering
\begin{tabular}{lrrrr}
\toprule
graph & bound $k$ & $|\Rel_k|$ & time (s) & time of \cref{alg:bucket} (s)\\
\midrule
$g_1$ & $\vs=11$ & 9\,500 & 0.62 & 0.22\\
$g_1$ & 13 & 55\,297 & 4.7 & \\
$g_1$ & 15 & 232\,170 & 28.6 & \\
$g_3$ & $\vs=15$ & 44\,952 & 3.7 & 1.45\\
$g_3$ & 17 & 258\,033 & 24.0 & \\
\bottomrule
\end{tabular}
\caption{Cost of overshooting $k$ with \cref{alg:kob} (Python, same
machine).}\label{tab:overshoot}
\end{table}

\paragraph{Comparison with MIP and MaxSAT.}
\Cref{tab:mip} compares the C implementation of \cref{alg:bucket} with the
MaxSAT model of \cref{sec:maxsat} on the same machine. On these graphs the
C implementation is between roughly 300 and 1500 times faster. On the
author's own machine the gap on $Q_{63}$ was larger still: 2637\,s for the
MaxSAT model against 0.89\,s for the C implementation, a ratio of about
3000. On those graphs
the MaxSAT model was in turn often hundreds of times faster than Coudert's
MIP solved with Gurobi~\cite{Gurobi}.

\begin{table}[ht]
\centering
\begin{tabular}{lrrrrrr}
\toprule
graph & $n$ & $\vs$ & MIP/Gurobi (s) & MaxSAT/RC2 (s) & \cref{alg:bucket}, C (s) & MaxSAT\,/\,C\\
\midrule
$g_1$ & 40 & 11 & \emph{[tbd]} & 11.9 & 0.008 & $\approx 1500$\\
$g_3$ & 45 & 15 & \emph{[tbd]} & 15.4 & 0.019 & $\approx 800$\\
$g_2$ & 60 & 11 & \emph{[tbd]} & 305 & 1.04 & $\approx 290$\\
$Q_{63}$ & 63 & 14 & \emph{[tbd]} & 198 & 0.43 & $\approx 460$\\
\bottomrule
\end{tabular}
\caption{Exact methods compared on one core of the machine described above.
The MaxSAT times are for RC2 with the CaDiCaL~1.9.5 backend (PySAT~1.9),
including model construction. The C times include reading the graph and
extracting and checking the ordering. The MIP column is to be filled in from
the author's Gurobi runs.}\label{tab:mip}
\end{table}

\paragraph{Comparison with branch and bound.}
We compared the C implementation with Sage's \texttt{vertex\_separation\_BAB}
(Cython; passagemath~10 distribution~\cite{Sage}), which implements the
greedy steps and prefix table of~\cite{CMN14} but not their reduction rules.
We used its default limits: prefixes of length at most 20, and at most
$10^6$ stored prefixes. The instances were the square-sum graphs
$Q_{30},\dots,Q_{63}$, four grids, 36 random graphs $G(n,p)$ with
$n\in\{30,40,50,60\}$ and average degree $3$, $4$ or $5$ (three seeds
each), and 9 random digraphs. BAB was timed in-process, with a cap of
120\,s; the C program was timed including process start-up. Whenever both
finished, the widths agreed. \Cref{tab:bab} shows the harder instances and
\cref{tab:babsum} a summary.

\begin{table}[ht]
\centering
\small
\begin{tabular}{lrrrrrrrr}
\toprule
& & & & \multicolumn{2}{c}{C, $U^*$} & \multicolumn{2}{c}{C, $U^{**}$} & BAB\\
\cmidrule(lr){5-6}\cmidrule(lr){7-8}
graph & $n$ & $m$ & $\vs$ & s & sets & s & sets & s\\
\midrule
$Q_{55}$ & 55 & 108 & 10 & 0.017 & 18\,126 & 0.009 & 8\,438 & 0.049\\
$Q_{57}$ & 57 & 114 & 11 & 0.029 & 34\,445 & 0.027 & 18\,009 & 0.169\\
$Q_{58}$ & 58 & 117 & 12 & 0.093 & 96\,416 & 0.060 & 46\,206 & 0.564\\
$Q_{59}$ & 59 & 120 & 13 & 0.279 & 249\,962 & 0.156 & 92\,992 & 2.080\\
$Q_{61}$ & 61 & 127 & 14 & 0.775 & 596\,416 & 0.330 & 188\,251 & 4.706\\
$Q_{63}$ & 63 & 135 & 14 & 0.422 & 364\,589 & 0.252 & 157\,639 & 2.713\\
grid $8\times8$ & 64 & 112 & 8 & 0.010 & 11\,045 & 0.005 & 1\,442 & 0.044\\
$G(50,5,1)$ & 50 & 119 & 14 & 0.180 & 131\,856 & 0.069 & 31\,490 & 1.309\\
$G(50,5,2)$ & 50 & 132 & 15 & 0.210 & 179\,617 & 0.091 & 42\,972 & 1.576\\
$G(60,3,0)$ & 60 & 104 & 11 & 1.780 & 919\,483 & 0.493 & 143\,089 & 12.2\\
$G(60,3,1)$ & 60 & 104 & 10 & 0.469 & 319\,193 & 0.168 & 67\,078 & 3.241\\
$G(60,4,0)$ & 60 & 125 & 14 & 5.165 & 3\,019\,891 & 1.205 & 317\,110 & $>120$\\
$G(60,4,1)$ & 60 & 128 & 13 & 1.545 & 998\,184 & 0.538 & 167\,512 & $>120$\\
$G(60,4,2)$ & 60 & 120 & 12 & 0.821 & 643\,444 & 0.318 & 146\,301 & 7.866\\
$G(60,5,0)$ & 60 & 158 & 16 & 0.753 & 548\,970 & 0.403 & 224\,144 & 5.378\\
$G(60,5,1)$ & 60 & 159 & 17 & 1.565 & 1\,194\,285 & 0.889 & 396\,736 & $>120$\\
$G(60,5,2)$ & 60 & 149 & 16 & 1.714 & 1\,200\,576 & 0.630 & 230\,728 & $>120$\\
\bottomrule
\end{tabular}
\caption{The C implementation of \cref{alg:bucket}, with Kobayashi's
expansion $U^*$ and with the stronger expansion $U^{**}$
(\cref{sec:greedy2}), against the branch and bound of~\cite{CMN14} in Sage
(cap 120\,s). $G(n,d,s)$ denotes a random graph with $n$ vertices,
average degree $d$ and seed $s$.}\label{tab:bab}
\end{table}

\begin{table}[ht]
\centering
\small
\begin{tabular}{lrrrrr}
\toprule
family & instances & BAB $>120$\,s & BAB / C ($U^*$) & BAB / C ($U^{**}$) & sets $U^*$ / $U^{**}$\\
\midrule
square-sum $Q_{30..63}$ & 12 & 0 & 5.8 & 9.5 & 2.4\\
grids & 4 & 0 & --- & --- & 6.1\\
$G(30,\cdot,\cdot)$ & 9 & 0 & --- & --- & 2.6\\
$G(40,\cdot,\cdot)$ & 9 & 0 & 7.4 & 11.1 & 2.9\\
$G(50,\cdot,\cdot)$ & 9 & 0 & 7.7 & 19.9 & 4.5\\
$G(60,\cdot,\cdot)$ & 9 & 4 & 7.6 & 21.8 & 5.1\\
random digraphs & 9 & 0 & --- & --- & 2.4\\
\bottomrule
\end{tabular}
\caption{Geometric means of time ratios, over the instances on which BAB
finished and the C program ($U^*$) took at least 10\,ms; shorter runs are
dominated by process start-up. The last column covers all instances.}\label{tab:babsum}
\end{table}

With its default limits BAB exceeded the cap on four $G(60,\cdot,\cdot)$
instances. Rerun with prefixes of any length and up to $10^7$ stored
prefixes, it solved them in 12.3, 43.2, 19.5 and 13.3\,s ($G(60,4,1)$,
$G(60,4,0)$, $G(60,5,1)$, $G(60,5,2)$). These times are 8--13 times the C
times with $U^*$, and 21--36 times those with $U^{**}$. With the default
limits, $G(60,4,1)$ still had not finished after 600\,s. So the timeouts
come from the cap on stored prefixes, not from the search itself. We did
not rerun the other instances with raised limits.

The two programs are much closer to each other than either is to the
MaxSAT model, and there is a structural reason. A branch and bound that
proves optimality must refute every prefix of width below the final upper
bound $\vs(D)$. With a memo on prefix sets, it therefore explores a family
of greedily closed sets comparable to $\Rel_{\vs(D)-1}$, the lower end of
the sandwich in \cref{thm:work}, which \cref{alg:bucket} also visits. (This
is a heuristic explanation; the two expansion rules and the depth-first
order make the families differ in detail.) The remaining differences are constants (bit-parallel C against
Cython with a Python set of frozensets), the stronger expansion, and BAB's
cap on stored prefixes. Once the cap is reached, sets deeper than 20 are
no longer memoised, and they can be re-explored many times.

\paragraph{Optimal orderings and secondary costs.}
\Cref{tab:orders} shows the complete DAG $\Gamma$ for square-sum graphs, and
the time to minimise a step-local proxy cost over it. The cost is the one
the author uses to tune a frontier dynamic program for counting Hamiltonian
paths. Writing $f$ for the new frontier size, $b$ for the number of edges
from $v$ into the frontier, $\epsilon$ for the number of edges inside the
frontier, $\rho\le1$ for a component correction, and $\ell$ for the number of
eliminations, a step contributes
$e=1.55^{b}(f/k)^2\rho\,(1+\epsilon/(4f))$,
$\kappa=\max(2^{-\ell},0.01)$ and $a=0.001f$. The Pareto fronts never held
more than one label on these graphs; on a random graph $g_1$ they reached
six. We checked our implementation of the cost against the author's code to
$10^{-9}$ on sampled orderings and on the optima. For comparison, the family
of \emph{all} optimal orderings, enumerated by their prefix sets, has
$5.36$M sets for $Q_{50}$, $8.97$M for $Q_{55}$ and $51.7$M for $Q_{57}$, and
did not fit in 7\,GB for $Q_{58}$.

\begin{table}[ht]
\centering
\small
\begin{tabular}{lrrrrrr}
\toprule
graph & $\vs$ & $|\Rel_{\vs}|$ & nodes of $\Gamma$ & orderings & C, complete (s) & optimise (s)\\
\midrule
$Q_{57}$ & 11 & 99\,463 & 13\,762 & $1.6\times10^{26}$ & 0.12 & 2\\
$Q_{58}$ & 12 & 278\,855 & 45\,056 & $3.6\times10^{28}$ & 0.38 & 9\\
$Q_{59}$ & 13 & 762\,813 & 151\,041 & $2.0\times10^{31}$ & 1.32 & 33\\
$Q_{61}$ & 14 & 1\,665\,113 & 263\,672 & $2.2\times10^{33}$ & 3.37 & 66\\
\bottomrule
\end{tabular}
\caption{The complete DAG $\Gamma$. The orderings column counts paths, with
one canonical order inside each expansion block. The optimisation of the
proxy cost was done in Python.}\label{tab:orders}
\end{table}

How much is lost by restricting to $\Gamma$? \Cref{tab:gammaloss} compares
the exact optimum over $\Gamma$ with the exact optimum over all optimal
orderings, for an additive cost and for the proxy cost. On these graphs the
loss is at most 1.5\%. Restricting further to the stronger expansion costs
much more: for $Q_{40}$ the proxy optimum rises from 3.59 to 4.74, and for
$Q_{45}$ from 1.45 to 2.43. These are empirical observations on small
graphs.

\begin{table}[ht]
\centering
\small
\begin{tabular}{lrrrrrr}
\toprule
& & \multicolumn{2}{c}{additive (V2)} & \multicolumn{2}{c}{proxy (V7)} & \\
\cmidrule(lr){3-4}\cmidrule(lr){5-6}
graph & $\vs$ & all & $\Gamma$ & all & $\Gamma$ & largest ratio\\
\midrule
$Q_{20}$ & 3 & 9.032 & 9.143 & 6.372 & 6.380 & 1.012\\
$Q_{22}$ & 3 & 15.599 & 15.710 & 9.510 & 9.519 & 1.007\\
$Q_{25}$ & 4 & 18.629 & 18.691 & 7.780 & 7.814 & 1.004\\
$Q_{28}$ & 5 & 20.730 & 20.770 & 5.615 & 5.698 & 1.015\\
\bottomrule
\end{tabular}
\caption{Exact optimum of two secondary costs over all pathwidth-optimal
orderings and over those of $\Gamma$. V2 is
$\sum_s 2^{b}(f/k)^2+0.001\sum_s f$.}\label{tab:gammaloss}
\end{table}

These comparisons should be read with the different scaling of the methods
in mind. The running time of \cref{alg:bucket} is governed by
$|\Rel_{\vs(D)}|\le 3^{\tau(D)}$. It is therefore at its best on graphs with
small vertex cover number, or more precisely few relevant sets, whatever the
value of $n$, and it can run out of memory on graphs with large $\tau$. The
performance of the MIP and MaxSAT approaches depends on quite different
structural features.

\FloatBarrier
\section{Conclusion}

Kobayashi's relevant-set dynamic program has a $k$-independent transition
cost. This turns the decision procedure into a single bottleneck
shortest-path computation. The resulting algorithm finds the exact pathwidth
and an optimal ordering in one pass, with the same $3^{\tau(D)}n^{O(1)}$
worst-case bound, and its work is sandwiched between one decision run just
below the optimum and one at the optimum. The only delicate point in
extracting the ordering is to emit the vertices of each expansion step in
the order in which they become subsumed.

The second greedy rule of Coudert, Mazauric and Nisse fits into the
expansion and reduces the work further. Completing the final level turns
the search into a compact representation of a very large family of optimal
orderings. Over that family, step-local secondary costs, including costs
with a multiplicative carry, can be minimised exactly, rather than searched
by local moves. On our graphs the restriction to this family cost at most
1.5\% against the optimum over all optimal orderings, and that larger family
quickly became too big to enumerate.

The method is not specific to pathwidth. It needs three things: a step
cost that depends only on the prefix set and the next vertex, a min--max
objective, and greedy rules that do not depend on the width bound. For a
submodular prefix cost the rules come from a single lemma, which explains
both published rules for pathwidth and gives exact cutwidth over at most
$2^{\tau}$ states. Treewidth satisfies the first two conditions but has
weak rules and much stronger dedicated methods. Bandwidth fails the first.
Linear rank-width, whose cut-rank function is submodular, is a natural next
candidate.

\begin{thebibliography}{99}
\bibitem{BFKKT12}
H.~L. Bodlaender, F.~V. Fomin, A.~M.~C.~A. Koster, D.~Kratsch, D.~M. Thilikos,
A note on exact algorithms for vertex ordering problems on graphs,
\emph{Theory Comput.\ Syst.} 50(2) (2012) 420--432.

\bibitem{CLTV13}
M.~Chapelle, M.~Liedloff, I.~Todinca, Y.~Villanger,
Treewidth and pathwidth parameterized by the vertex cover number,
in: \emph{Proc.\ WADS 2013}, LNCS 8037, 2013, pp.~232--243.

\bibitem{CMN14}
D.~Coudert, D.~Mazauric, N.~Nisse,
Experimental evaluation of a branch and bound algorithm for computing
pathwidth,
Research report RR-8470, Inria, 2014, hal-00943549.

\bibitem{Cou16}
D.~Coudert,
A note on integer linear programming formulations for linear ordering
problems on graphs,
Research report hal-01271838, Inria, 2016.

\bibitem{CouCode}
D.~Coudert,
Linear ordering problems on graphs (code),
\url{https://www-sop.inria.fr/members/David.Coudert/code/graph-linear-ordering.shtml}.

\bibitem{CLPPS14}
M.~Cygan, D.~Lokshtanov, M.~Pilipczuk, M.~Pilipczuk, S.~Saurabh,
On cutwidth parameterized by vertex cover,
\emph{Algorithmica} 68(4) (2014) 940--953.

\bibitem{CP12}
M.~Cygan, M.~Pilipczuk,
Even faster exact bandwidth,
\emph{ACM Trans.\ Algorithms} 8(1) (2012), article~8; arXiv:0902.1661.

\bibitem{BBNNPR13}
T.~Biedl, T.~Bl\"asius, B.~Niedermann, M.~N\"ollenburg, R.~Prutkin, I.~Rutter,
Using ILP/SAT to determine pathwidth, visibility representations, and other
grid-based graph drawings, arXiv:1308.6778, 2013.

\bibitem{Dij59}
E.~W. Dijkstra,
A note on two problems in connexion with graphs,
\emph{Numer.\ Math.} 1 (1959) 269--271.

\bibitem{Gurobi}
Gurobi Optimization, LLC,
Gurobi Optimizer Reference Manual, \url{https://www.gurobi.com}.

\bibitem{RC2}
A.~Ignatiev, A.~Morgado, J.~Marques-Silva,
RC2: an efficient MaxSAT solver,
\emph{J.\ Satisf.\ Boolean Model.\ Comput.} 11(1) (2019) 53--64.

\bibitem{PySAT}
A.~Ignatiev, A.~Morgado, J.~Marques-Silva,
PySAT: a Python toolkit for prototyping with SAT oracles,
in: \emph{Proc.\ SAT 2018}, LNCS 10929, 2018, pp.~428--437.

\bibitem{GS84}
E.~M. Gurari, I.~H. Sudborough,
Improved dynamic programming algorithms for bandwidth minimization and the
MinCut linear arrangement problem,
\emph{J.\ Algorithms} 5(4) (1984) 531--546.

\bibitem{KF79}
T.~Kashiwabara, T.~Fujisawa,
NP-completeness of the problem of finding a minimum-clique-number interval
graph containing a given graph as a subgraph,
in: \emph{Proc.\ ISCAS 1979}, pp.~657--660.

\bibitem{Kin92}
N.~G. Kinnersley,
The vertex separation number of a graph equals its path-width,
\emph{Inform.\ Process.\ Lett.} 42(6) (1992) 345--350.

\bibitem{KKKTT12}
K.~Kitsunai, Y.~Kobayashi, K.~Komuro, H.~Tamaki, T.~Tano,
Computing directed pathwidth in $O(1.89^n)$ time,
in: \emph{Proc.\ IPEC 2012}, LNCS 7535, 2012, pp.~182--193.

\bibitem{Kob15}
Y.~Kobayashi,
Computing the pathwidth of directed graphs with small vertex cover,
\emph{Inform.\ Process.\ Lett.} 115 (2015) 310--312.

\bibitem{Knu77}
D.~E. Knuth,
A generalization of Dijkstra's algorithm,
\emph{Inform.\ Process.\ Lett.} 6(1) (1977) 1--5.

\bibitem{Mal18}
S.~Mallach,
Linear ordering based MIP formulations for the vertex separation or
pathwidth problem,
\emph{J.\ Discrete Algorithms} 52--53 (2018) 156--168.

\bibitem{Pol60}
M.~Pollack,
The maximum capacity through a network,
\emph{Oper.\ Res.} 8(5) (1960) 733--736.

\bibitem{Sax80}
J.~B. Saxe,
Dynamic programming algorithms for recognizing small-bandwidth graphs in
polynomial time,
\emph{SIAM J.\ Algebraic Discrete Methods} 1(4) (1980) 363--369.

\bibitem{Sage}
The Sage Developers, SageMath; modularised as passagemath,
\url{https://www.sagemath.org}, \url{https://github.com/passagemath/passagemath};
module \texttt{sage.graphs.graph\_decompositions.vertex\_separation}.

\bibitem{SV09}
K.~Suchan, Y.~Villanger,
Computing pathwidth faster than $2^n$,
in: \emph{Proc.\ IWPEC 2009}, LNCS 5917, 2009, pp.~324--335.

\bibitem{Tam11}
H.~Tamaki,
A polynomial time algorithm for bounded directed pathwidth,
in: \emph{Proc.\ WG 2011}, LNCS 6986, 2011, pp.~331--342.

\bibitem{Tam19}
H.~Tamaki,
Positive-instance driven dynamic programming for treewidth,
\emph{J.\ Comb.\ Optim.} 37(4) (2019) 1283--1311; conference version in
\emph{Proc.\ ESA 2017}, LIPIcs 87, 68:1--68:13.

\bibitem{YC08}
B.~Yang, Y.~Cao,
Digraph searching, directed vertex separation and directed pathwidth,
\emph{Discrete Appl.\ Math.} 156(10) (2008) 1822--1837.
\end{thebibliography}

\end{document}
